\documentclass{homework}
% \usepackage{graphicx} % Required for inserting images
\usepackage{amsmath}
\usepackage{gensymb}
\usepackage{subfig}
\usepackage{float}

\class{ECEN 5448}
\title{Midterm 2}
\author{Sean Jennings}
\date{December 5, 2023}

\newcommand{\obs}{\mathcal{O}}

\begin{document}
\maketitle

\section{Problem 1}
To construct a quadratic Lyaponov function, we solve the Lyaponov equation:
\begin{equation*}
    PA + A^T P = -Q
\end{equation*}
where $A$ is the state dynamics matrix, $P$ is the symmtric matrix for the Lyaponov function,
and $Q$ is some positive definite matrix. Let $Q$ be the identity, and $P_{ij}$ denote
the $ij$th entry of $P$. Then,
\begin{align*}
    \mat{
        P_{11} & P_{12} \\
        P_{12} & P_{22}
    } \mat{
        -1 & 5 \\
        1 & -10 \\
    } + \mat{
        -1 & 1 \\
        5 & -10
    } \mat{
        P_{11} & P_{12} \\
        P_{12} & P_{22}
    } &= \mat{
        -1 & 0 \\
        0 & -1 \\
    } \\
    \mat{
        -P_{11} + P_{12} & 5 P_{11} - 10 P_{12} \\
        -P_{12} + P_{22} & 5 P_{12} - 10 P_{22} \\
    } + \mat{
        -P_{11} + P_{12} & -P_{12} + P_{22} \\
        5 P_{11} - 10 P_{12} & 5 P_{12} - 10 P_{22} \\
    } &= \mat{
        -1 & 0 \\
        0 & -1 \\
    }
\end{align*}
yields the following system of equations:
\begin{align*}
    2(-P_{11} + P_{21}) &= -1 \\
    5 P_{11} - 11 P_{12} + P_{22} &= -1 \\
    10 P_{12} - 20 P_{12} &= -1.
\end{align*}
Converting this to the matrix form $Ax = b$ where $x = [P_{11}, P_{12}, P_{22}]^T$,
and solving via matrix inversion (with a computer) yields the solution:
\begin{equation*}
    P = \frac{1}{22} \mat{
        21.2 & 10.2 \\
        10.2 & 6.2 \\
    }
\end{equation*}

We can check the eigenvalues of $P$ (again plugging into Python), and find
\begin{align*}
    \lambda_1(P) &\approx 1.198 \\
    \lambda_2(P) &\approx 0.0472
\end{align*}
so $P$ is a positive definite matrix and $\epsilon(x)$ is a positive definite
function on $\R^2$. This also implies that $\epsilon$ is radially unbounded. By satisfying
the Lyaponov equation, we have
\begin{align*}
    \dot{\epsilon}(x) &= x^T P \dot{x} + \dot{x}^T P x  \\
        &= x^T (PA + A^T P) x \\
        &= - x^T Q x \\
        &= - \norm{x}.
\end{align*}
So, $\dot{\epsilon}(x)$ is a negative definite function on the domain of $\R^2$.
Thus, by Lyaponov's second method, the conditions for a globally asymptotically
stable equilibrium are satisfied.

\section{Problem 2}

\begin{letterenum}
\item We will define $x\in \R^2$ to be
\begin{equation*}
    x := \mat{
        x_1 \\ \tilde{x}_2
    }.
\end{equation*}
We can then rewrite the system in matrix form:
\begin{equation*}
    \dot{x} = \mat{
        -K & K \\
        K & -K
    } x + \mat{1 \\ 0} u
\end{equation*}
or with $K=1$:
\begin{equation*}
    \dot{x} = \mat{
        -1 & 1 \\
        1 & -1
    } x + \mat{1 \\ 0} u
\end{equation*}

\newcommand{\contr}{\mathcal{C}}
Since our system is LTI, we can look at the the controllability matrix, $\contr(A, B)$
to determine controllability. We compute the controllability matrix to be:
\begin{equation*}
    \contr(A, B) = \mat{B & AB} = \mat{
        1 & -1 \\
        0 & 1 \\
    }.
\end{equation*}
Clearly, $\contr(A, B)$ is full rank and therefore our system is controllable.

To find the matrix in controllable canonical form (CCF), we first determine the characteristic 
polynomial:
\begin{align*}
    I\lambda - A &= \mat{
        \lambda + 1 & -1 \\
        - 1 & \lambda + 1 \\
    } \\
    \det (I\lambda - A) &= (\lambda + 1)^2 - 1 = \lambda^2 + 2 \lambda
\end{align*}

The CCF system $(A', B')$ is determined uniquely by the characteristic polynomial to be
\begin{equation*}
    A' = \mat{
        0 & 1 \\
        0 & -2
    } \qquad B' = \mat{
        0 \\ 1
    }
\end{equation*}
The transformation $P$ which brings the system into CCF is then given by:
\begin{align*}
    P &= \contr(A', B') \contr(A, B)^{-1} \\
        &= \mat{
            0 & 1 \\
            1 & -2 \\
        } \mat{
            1 & 1 \\
            0 & 1
        } \\
        &= \mat{
            0 & 1 \\
            1 & -1 \\
        }
\end{align*}

\item To find the control input which drives the system from state
$x(0) = [3, 1]^T$ to $x(1) = [4, 3]^T$, we look at the minimum
effort controller:
\begin{equation*}
    u(\tau) = B^T \Phi^T(1, \tau) v
\end{equation*}
where $v\in\R^2$ is the vector such that
\begin{equation*}
    x(1) - \Phi(1, 0) x(0) = W_R(0, 1) v
\end{equation*}
with $\Phi(t, t_0): \R \times \R \to \R^{2\times 2}$ representing
the state transition matrix and $W_R(0, 1)$ representing
the reachability grammian on the interval $[t_0, t_1] =[0, 1]$.

We must therefore compute the state transition matrix and reachability
grammian. First, the state transition matrix is found by 
diagonalization of the $A$ matrix:
\begin{equation*}
    A = \mat{
        -1 & 1 \\
        1 & -1
    } = PDP^{-1} \qquad P = \mat{
        1 & 1 \\
        -1 & 1
    } \qquad D = \mat{
        -2 & 0 \\
        0 & 0
    }
\end{equation*}
The state transition matrix is then given by:
\begin{align*}
    \Phi(t, t_0) = e^{A(t - t_0)} &= P e^{D(t-t_0)} P^{-1} \\
        &= \mat{
            1 & 1 \\
            -1 & 1
        } \mat{
            e^{-2(t-t_0)} & 0 \\
            0 & 1
        } P^{-1} \\
        &= \mat{
            e^{-2(t-t_0)} & 1 \\
            -e^{-2(t-t_0)} & 1 \\
        } \frac{1}{2} \mat{
            1 & -1 \\
            1 & 1 \\
        } \\
\end{align*}
and a final expression for the state transition matrix of
\begin{equation}
        \Phi(t, t_0) = \frac{1}{2} \mat{
            1 + e^{-2(t - t_0)} & 1 - e^{-2(t - t_0)} \\
            1 - e^{-2(t - t_0)} & 1 + e^{-2(t - t_0)}
        }
        \label{eq:state-transition}
\end{equation}
is obtained.

Then, we can compute the reachability grammian. We will define
the variable $a := e^{-2(t - t_0)}$ to simplify the exposition:
\begin{align*}
    W_R(0, 1) &= \int_{0}^{1} \Phi(1, \tau) B B^T \Phi^T(1, \tau) d\tau \\
        &= \int_{0}^{1} \Phi(1, \tau) \mat{1 \\ 0} \mat{1 & 0} \Phi^T(1, \tau) d\tau \\
        &= \int_{0}^{1} \frac{1}{2} \mat{
            1 + a & 1 - a \\
            1 - a & 1 + a
        } \mat{1 & 0 \\ 0 & 0} \Phi^T(1, \tau) d\tau \\
        &= \frac{1}{4} \int_{0}^{1} \mat{
            1 + a & 0 \\
            1 - a & 0 
        } \mat{
            1 + a & 1 - a \\
            1 - a & 1 + a \\
        } d\tau \\
        &= \frac{1}{4} \int_{0}^{1} \mat{
            (1 + a)^2 & (1 + a)(1 - a) \\
            (1 + a)(1 - a) & (1 - a)^2 \\
        } d\tau.
\end{align*}
We next define the variables $b,c,d\in\R$ to be the integral of the following matrix components:
\begin{align*}
    b :&= \int_{0}^{1} (1 + a)^2 d\tau  \\
        &= \int_{0}^{1} 1 + 2 e^-2(1 - \tau) + e^{-4(1 - \tau)} d\tau \\
        &= \tau + e^{-2(1-\tau)} + \frac{1}{4} e^{-4(1-\tau)} \eval_{\tau = 0}^1 \\
        &= (1 + 1 + \frac{1}{4}) - (e^{-2} + e^{-4}) \\
        &= \frac{9}{4} - e^{-2} - e^{-4} \\
    c :&= \int_{0}^{1} (1 + a)(1 - a) d\tau  \\
        &= \int_{0}^{1} 1 - e^{-4(1 - \tau)} d\tau \\
        &= \tau - \frac{1}{4} e^{-4(1-\tau)} \eval_{\tau = 0}^1 \\
        &= (1 - \frac{1}{4}) - (\frac{1}{4} e^{-4}) \\
        &= \frac{3}{4} - \frac{1}{4} e^{-4} \\
    d :&= \int_{0}^{1} (1 - a)^2 d\tau  \\
        &= \int_{0}^{1} 1 - 2 e^-2(1 - \tau) + e^{-4(1 - \tau)} d\tau \\
        &= \tau - e^{-2(1-\tau)} + \frac{1}{4} e^{-4(1-\tau)} \eval_{\tau = 0}^1 \\
        &= (1 - 1 + \frac{1}{4}) - (-e^{-2} + e^{-4}) \\
        &= \frac{1}{4} + e^{-2} - e^{-4}
\end{align*}
The grammian on the desired interval is then given by:
\begin{equation*}
    W_R(0, 1) = \frac{1}{4} \mat{
        b & c \\
        c & d \\
    }
\end{equation*}
Define $\Delta$ to be the determinant of the grammian:
\begin{equation*}
    \Delta := \text{det}(W_R(0, 1)) = \frac{1}{16} (bd - c^2)
\end{equation*}
The inverse of the grammian is then given by:
\begin{equation*}
    W_R(0, 1)^{-1} = \frac{1}{4\Delta} \mat{
        d & -c \\
        -c & b \\
    }
\end{equation*}
Finally, we can find an expression for the control input:
\begin{align*}
    u(t) &= B^T \Phi^T(1, \tau) W_R(0, 1)^{-1} \biggl[
        x(1) - \Phi(1, 0) x(0)
    \biggr] \\
    &= \mat{1 & 0} \mat{
        1 + a & 1 - a \\
        1 - a & 1 + a \\
    } \frac{1}{4\Delta} \mat{
        d & -c \\
        -c & b
    } \biggl(
        \mat{4 \\ 3} - \mat{
            1 + e^{-2} & 1 - e^{-2} \\
            1 - e^{-2} & 1 + e^{-2} \\
        } \mat{3 \\ 1}
    \biggr) \\
    &= \frac{1}{4\Delta} \mat{1 + a & 1 - a} \mat{
        d & -c \\
        -c & b
    } \mat{
        -2 e^{-2} \\
        2e^{-2} - 1
    } \\
    &= \frac{1}{4\Delta} \mat{
        (d - c) + (d + c) a & (b - c) - (b + c) a
    } \mat{
        -2e^{-2} \\
        2e^{-2} - 1
    }  \\
    &= \frac{1}{4\Delta} \biggl(
        2e^{-2}(b - d) - b + c + [ b + c - 2e^{-2}(b + 2c + d)]a
    \biggr) \\
    &\approx -14.62 + 29.15 a \\
    &= -14.62 + 29.15 e^{-2(1 - t)}
\end{align*}

\end{letterenum}

\section{Problem 3}

\begin{letterenum}

\item To show observability of an LTI system, we can check the nullspace
of the observability matrix. If the observability matrix $\obs(A, C)$ is
full column rank, then the null space is $\{0\}$ and the system is 
observable. $\obs(A, C)$ is given by
\begin{equation*}
    \obs(A, C) = \mat{
        C \\ CA
    } = \mat{
        1 & 0 \\
        -1 & 1 \\
    }.
\end{equation*}
We see that the columns of $\obs(A, C)$ are linearly independent and
therefore the nullspace is trivial and the system is observable.

\item The transfer function of the system is given by:
\begin{equation*}
    G(s) = C(sI - A)^{-1} B.
\end{equation*}
So we can first compute $(sI - A)^{-1}$:
\begin{equation*}
    (sI - A)^{-1} = \mat{
        s + 1 & - 1 \\
        -1 & s + 1 \\
    } = \frac{1}{s(s+2)} \mat{
        s + 1 & 1 \\
        1 & s + 1 \\
    }
\end{equation*}
and therefore the transfer function is:
\begin{align*}
    G(s) &= C (sI - A)^{-1} B \\
        &= \mat{1 & 0} \biggl(
            \frac{1}{s(s+2)} \mat{
                s + 1 & 1 \\
                1 & s + 1 \\
            }
        \biggr) \mat{
            1 \\ 0
        } \\
        &= \frac{s + 1}{s(s+2)}
\end{align*}

\item To reconstruct the states on the given interval, it suffices to 
determine the initial condition $x(0)$, since the states are given
by the LTI system solution
\begin{equation}
    x(t) = \Phi(t, 0) x(0) + \int_{0}^{t} \Phi(t, \tau) B u(\tau) d\tau.
    \label{eq:system-solution}
\end{equation}

To find the initial state, we note that $y(t) = Cx(t)$. We can substitute the
previous expression and subtract by the known integral control input term, to
define $\tilde{y}$, the control output compensating for the input:
\begin{align}
    \tilde{y}(t) :&= y(t) - C \int_{0}^{t} \Phi(t, \tau) B u(\tau) d\tau \label{eq:tilde-y}\\
        &= C \Phi(t, 0) x(0)
\end{align}
Multiplying both sides by $\Phi^T(t, 0) C^T$ and integrating over the interval $[0, t]$ relates
$\tilde{y}$ to the observability grammian $M_O(0, t)$:
\begin{equation*}
    \int_{0}^{t} \Phi^T(\tau, 0) C^T \tilde{y}(\tau) d\tau = M_O(0, t) x(0)
\end{equation*}
Since the system is observable, $M_O(0, t)$ is invertible over all time intevals.
Thus, the initial condition can be derived from the following expression:
\begin{equation}
    x(0) = M_O(0, 1)^{-1} \int_{0}^{1}\Phi^T(\tau, 0) C^T \tilde{y}(\tau) d\tau
    \label{eq:init-cond}
\end{equation}

We must therefore evaluate the grammian and the integral involving $\tilde{y}$ to compute the
output. We start with the definition for $\tilde{y}$ in equation (\ref{eq:tilde-y}), taking
$u(\tau) = 1$ and noting that left multiplication of $\Phi(t, 0)$ by $C=\mat{1 & 0}$ along with right multiplication
by $B = \mat{1 & 0}^T$ takes its top left element:
\begin{align*}
    \tilde{y}(t) &= y(t) - \int_{0}^{t} C\Phi(t, \tau) B d\tau \\
        &= y(t) - \int_{0}^{t} 1 + e^{-2(t - \tau)} d\tau \\
        &= y(t) - \biggl[
            \tau + \frac{1}{2} (e^{-2(t-\tau)})
        \biggr]\eval_{\tau=0}^t \\
        &= y(t) - (t + \frac{1}{2} - \frac{1}{2}e^{-2t})
\end{align*}

A simplification of terms in $y(t)$ (omitted for the sake of time) yields
\begin{equation*}
    y(t) = \frac{5}{4} + \frac{1}{2} t + \frac{3}{4} e^{-2t}
\end{equation*}
resulting in a final expression for $\tilde{y}$:
\begin{align*}
    \tilde{y}(t) &= \frac{3}{4} + \frac{1}{2}t + \frac{3}{4} e^{-2t} \\
        &= \frac{1}{4}(3 - 2t + 5e^{-2t})
\end{align*}
The integral term is then found (again omitting tedious computation):
\begin{equation*}
    \int_{0}^{1} \Phi^T(\tau, 0) C^T \tilde{y}(\tau) d\tau =
        \frac{1}{32} \mat{
            27 - 10 e^{-2} - 5 e^{-4} \\
            9 - 10 e^{-2} + 5 e^{-4} \\
        }
\end{equation*}
The observability grammian is computed as follows:
\begin{align*}
    M_O(0, 1) &= \int_{0}^{1} \Phi^T(\tau, 0) C^T C \Phi(\tau, 0) d\tau \\
        &= \int_{0}^{1} \frac{1}{2} \mat{
            1 + e^{-2\tau} & 1 - e^{-2\tau} \\
            1 - e^{-2\tau} & 1 + e^{-2\tau} \\
        } \mat{
            1 & 0 \\
            0 & 0
        } \Phi(\tau, 0) d\tau \\
        &= \frac{1}{4} \int_{0}^{1} \mat{
            1 + 2e^{-2\tau} + e^{-4\tau} & 1 - e^{-4\tau} \\
            1 - e^{-4\tau} & 1 - 2e^{-2\tau} + e^{-4\tau}
        } \\
        &= \frac{1}{4} \mat{
            \tau - e^{-2\tau} - \frac{1}{4} e^{-4\tau}  & \tau + \frac{1}{4}e^{-4\tau} \\
            \tau + \frac{1}{4}e^{-4\tau} & \tau + e^{-2\tau} - \frac{1}{4} e^{-4\tau}
        } \\
        &= \frac{1}{4} \mat{
            \frac{9}{4} - e^{-2} - \frac{1}{4}e^{-4} & \frac{3}{4} + \frac{1}{4} e^{-4} \\
            \frac{3}{4} + \frac{1}{4} e^{-4} & \frac{1}{4} + e^{-2} - \frac{1}{4} e^{-4}
        }
\end{align*}

We now end our symbolic endeavours and evaluate $x(0)$ via equation (\ref{eq:init-cond}) with
the aid of a computer:
\begin{equation*}
    x(0) \approx \mat{
        2.078 \\
        -1.5782
    }
\end{equation*}
The state as a function of time is therefore given by equation (\ref{eq:system-solution}), with
state transition matrix given by (\ref{eq:state-transition}) and $x(0)$ given by the previous
equation.







\end{letterenum}

\section{Problem 4}

\begin{letterenum}
\item We recompute the observability matrix with $C = \mat{1 & -1}$:
\begin{equation*}
    \obs (A, C) = \mat{
        1 & -1 \\
        -2 & 2 \\
    }
\end{equation*}
and we can see that the nullspace is now nontrivial. The unobservable
subspace is equal to the nullspace of the observability matrix
for LTI systems and is therefore given by:
\begin{align*}
    \mathcal{U}0(t_0, t_1) &= \mathcal{N}(\obs(A, C)) \\
        &= \biggl\{
            x \in \R^2: \obs(A, C)x = 0
        \biggr\} \\
        &= \biggl\{
            x\in\R^2: x = \mat{1 \\ 1} \alpha, \alpha \in \R
        \biggr\}
\end{align*}
This unobservable subspace represents the set of $[x_1, \tilde{x}_2]$
with $x_1 = \tilde{x}_2$. That is, $x_1$ and $x_2$ are exactly $d$
distance apart. Given the trucks are $d$ apart, the position
of each truck is determined by $x_1$ alone. However, since we only
have relative measurements between the trucks, the absolute
position in space cannot be determined from the measurements.
We can determine positions only in relative terms. Any change 
in position which affects both trucks equally cannot be observed.

\item To find the transfer function, we use the value of $(sI - A)^{-1}$
computed in part 3b, but consider the new $C$ matrix:
\begin{align*}
    G(s) &= C(sI - A)^{-1} B \\
        &= \mat{1 & -1} \biggl(
            \frac{1}{s(s+2)} \mat{
                s + 1 & 1 \\
                1 & s + 1
            }
        \biggr) \mat{1 \\ 0} \\
        &= \frac{1}{s(s+2)} \mat{1 & -1} \mat{s + 1 \\ 1} \\
        &= \frac{s}{s(s+2)} = \frac{1}{s+2}
\end{align*}

We note that since one of the polynomial terms cancelled out, the
minimum realization is now just a first order system, given by:
\begin{align*}
    \dot{x} &= -2x + u \\
    y &= x
\end{align*}
with $x,u,y \in \R$.

\item In Homework 4, we proved that two initial conditional are output
equivalent iff their difference belongs to the nullspace of the
observability grammian. For an LTI system, we have
\begin{equation*}
    \mathcal{N}(\obs(A, C)) = \mathcal{N}(M_0([t_0, t_1]))
\end{equation*}
where $M_0([t_0, t_1])$ is the observability grammian on the interval
$[t_0, t_1]$. So it suffices to check that the difference between
the given initial conditions $x(0)\in\R^2$ and $x'(0)\in\R^2$ belongs to the
nullspace of the observability matrix:
\begin{align*}
    \obs(A,C) (x'(0) - x(0)) &= \mat{
        1 & -1 \\
        -2 & 2
    } \biggl(
        \mat{3 \\ 1} - \mat{2 \\ 0}
    \biggr) \\
    &= \mat{
        1 & -1 \\
        -2 & 2
    } \mat{1 \\ 1} \\
    &= \mat{0 \\ 0}
\end{align*}
\end{letterenum}









\end{document}
